\chapter{技术路线争议与开放研究问题}
\label{chap:technical-debates}

\section{争议必须改写成可比较实验}

“端到端还是模块化”“数据还是模型”若不固定任务、预算、失败代价和接口，就只是口号。一个有效争议至少给出竞争命题、共同评测集、控制变量、失败分层和否证阈值。本章不预测赢家，而是列出什么证据会迫使我们修改判断。

\section{端到端与模块化：学习边界放在哪里}

RT-2 与 OpenVLA 证明视觉语言表征可以直接进入动作预测\cite{brohan2023rt2,kim2024openvla}。SayCan 和行为树则把高层选择与可执行技能、显式控制流分开\cite{ahn2022saycan,colledanchise2018behaviortrees}。端到端路线减少手写接口，有机会利用联合数据；模块化路线便于局部验证、替换和故障隔离。二者并非二元：实际系统常把感知和策略学习化，同时保留低层控制、安全监控与恢复状态机。

现有论文很少在相同硬件、数据、算力和安全约束下比较两种边界。可证伪条件是：若端到端系统在未见对象/干扰上持续降低人工介入，同时保持最坏情况停止距离和故障定位能力，强模块分解的必要性下降；若小分布漂移反复产生不可诊断失效，且模块化替换能以较低数据成本恢复，则“统一模型自然更稳健”被否证。Simplex 与控制屏障函数说明学习策略外仍可放置可验证安全层\cite{sha2001simplex,ames2017cbf}。

\section{通用与垂类：复用收益能否覆盖验证成本}

Octo、OpenVLA 和 Open X-Embodiment 探索跨任务、跨数据源训练\cite{octo2024,kim2024openvla,openx2024}。ACT 与 Diffusion Policy 则在受控任务上展示动作建模\cite{zhao2023act,chi2023diffusionpolicy}。通用模型的价值不是参数更多，而是新任务所需数据、工程和停机验证更少；垂类系统的优势是分布、工具和验收边界清楚。

缺失实验应同时报告零样本成功、少样本适配曲线、旧任务回退、推理资源和安全再验证工时。若通用预训练在多个本体上把达到同一可靠性所需现场数据和工程时长稳定降低，且不增加关键失败，通用路线成立；若每个工位仍需大量重采、重新标定和安全审查，所谓通用性只是初始化。

\section{世界模型：显式预测何时值得}

PlaNet 与 DreamerV3 展示了在潜空间预测并学习价值的路线\cite{hafner2019planet,hafner2023dreamerv3}。显式世界模型适合反事实、长时回报和数据复用，但机器人接触、多主体和部分可观测环境会迅速放大模型偏差。反应式策略可避开长滚动误差，却难以解释长期资源和前置条件。

公平实验需固定真实交互预算，让世界模型与无模型/短视策略面对相同遮挡、接触和环境变化，并按预测时域报告校准误差。若更长想象时域没有改善闭环任务，或性能增益可由更大行为数据解释，则“预测是关键机制”证据不足；若在稀缺真机数据下显著减少危险试错并改善长时恢复，显式模型价值成立。

\section{动作表示：没有脱离控制频率的最佳 token}

动作块 ACT、连续扩散、流匹配与离散 token 分别改变多峰建模、平滑性、采样时延和错误传播\cite{zhao2023act,chi2023diffusionpolicy,lipman2023flowmatching,pertsch2025fast}。第 22--24 章已说明动作空间还受坐标系、频率、低层控制器和数据采样影响。

需要在相同观测、骨干、数据和低层接口下比较动作表示，报告端到端时延、接触峰值、重规划频率和恢复。若某表示只在离线误差上占优却增加接触冲击或闭环滞后，不能称为更优；跨任务稳定收益且对频率/本体扰动不敏感，才支持其通用优势。

\section{数据与模型规模：有效覆盖不是样本总数}

Open X-Embodiment、DROID、$\pi_0$ 和 GR00T 把跨机构数据、规模化采集与基础模型连接起来\cite{openx2024,khazatsky2024droid,black2024pi0,nvidia2025groot}。但 episode 数不能表示接触状态、失败、恢复、对象和本体覆盖；混合数据还会引入动作语义与标定不一致。

规模实验应绘制数据量、模型量、覆盖度和真实成功率的联合曲线，并保留来源/任务/本体外测试。若增加近重复数据或参数在分布外不再改善，瓶颈已转向覆盖、接口或控制；若控制覆盖后仍有稳定幂律收益，继续扩展规模才有依据。

\begin{table}[H]
\begingroup\hbadness=10000
\centering\small
\caption{五组争议的缺失实验与否证信号}
\label{tab:debates-falsification}
\begin{tabularx}{\textwidth}{p{0.15\textwidth}YYp{0.27\textwidth}}
\toprule
争议 & 共同实验 & 关键指标 & 否证信号 \\
\midrule
端到端/模块化 & 同硬件数据与安全约束 & 介入、诊断、停止边界 & 成功率高但失效不可隔离 \\
通用/垂类 & 多任务适配与回退 & 数据/工时、旧任务保持 & 每任务仍需完整重做 \\
世界模型 & 同真机交互预算 & 校准、长时回报、危险试错 & 长预测不改善闭环 \\
动作表示 & 同骨干与控制接口 & 时延、冲击、恢复 & 离线优而闭环差 \\
数据/模型规模 & 覆盖受控的缩放实验 & 分布外成功、边际收益 & 近重复扩展无收益 \\
\bottomrule
\end{tabularx}
\endgroup
\end{table}

\section{五个开放问题的可测门槛}

灵巧操作需要跨物体、材料和接触模式报告损伤与恢复；长时程自主需要在小时/天尺度报告无人工任务时长和故障链；跨本体迁移需要把适配数据、映射误差和低层重整定分开；端侧模型需要在功耗、热稳态和最坏时延下比较能力；持续学习需要证明新数据上线后旧能力、安全案例和可回滚性不退化。

\section{知识小结与综述判断}

路线优劣由任务分布、闭环频率、数据预算、失效代价和验证责任共同决定。综述判断是：短期最可信的系统仍会混合学习与显式工程边界；长期是否收敛到更统一模型，取决于其能否在真实暴露量下同时降低人工介入、适配成本和不可诊断风险。
